% !TeX root=main.tex
\chapter{پروژه‌های مهم منتخب}

\section*{اپلیکیشن فاکتور}
«فاکتور» یک اپلیکیشن حسابداری و مدیریت مالی است که توسط اینجانب طراحی و توسعه داده شده و تاکنون بیش از ۱۲۰{,}000 نصب فعال در کافه‌بازار داشته است. 
این اپلیکیشن بر پایه فریم‌ورک \lr{Flutter} توسعه یافته و از معماری‌های مدرن (Clean Architecture, Riverpod) بهره می‌برد. 
چالش اصلی این پروژه مدیریت بهینه داده‌ها، پشتیبانی از حالت آفلاین، و مقیاس‌پذیری برای تعداد بالای کاربران بود که با طراحی دقیق و استفاده از فناوری‌هایی مانند \lr{ObjectBox} و \lr{Retrofit} حل شد. 

\section*{پایان‌نامه کارشناسی‌ارشد}
پایان‌نامه من بر محور ترکیب اپلیکیشن‌های اندرویدی با الگوریتم‌های هوش مصنوعی و شبکه‌های عصبی متمرکز بوده است. 
در این پروژه از معماری‌های مختلف یادگیری عمیق (\lr{CNN}، \lr{RNN}، \lr{Transfer Learning}) استفاده شد تا قابلیت هوشمندسازی اپلیکیشن‌ها در حوزه‌های مختلف مانند دسته‌بندی تصاویر و تحلیل داده‌های متنی بررسی شود. 
این تجربه موجب افزایش تسلط من بر چارچوب‌های یادگیری عمیق و کاربرد آن‌ها در اپلیکیشن‌های واقعی شد. 

\section*{پروژه‌های پردازش تصویر دانشگاهی}
در طول دوران تحصیل، روی چندین پروژه پردازش تصویر کار کرده‌ام. 
این پروژه‌ها شامل شناسایی الگوها در تصاویر میکروسکوپی، تشخیص اشیاء و آنالیز داده‌های تصویری در محیط‌های آزمایشگاهی بوده‌اند. 
برای این کار از کتابخانه‌هایی مانند \lr{OpenCV} و شبکه‌های عصبی کانولوشنی استفاده کردم. 
این پروژه‌ها زیربنای علاقه‌مندی من به کاربردهای نوآورانه هوش مصنوعی در حوزه معدن شدند. 

\section*{پکیج اختصاصی PrimaryKit}
برای بهبود سرعت توسعه در پروژه‌های مختلف، یک پکیج اختصاصی به نام \lr{PrimaryKit} طراحی و توسعه داده‌ام. 
این پکیج شامل کامپوننت‌های مشترک (مانند دکمه‌ها و فیلدهای متنی) بوده و در پروژه‌های متعدد \lr{Flutter} استفاده شده است. 
این اقدام علاوه بر صرفه‌جویی در زمان توسعه، موجب افزایش هماهنگی در کدها و استانداردسازی پروژه‌ها شد.
